\documentclass[11pt,a4paper]{article}
% Ortak önbölüm. pdflatex ve xelatex ile derlenir.
\usepackage{iftex}
\ifPDFTeX
  \usepackage[T1]{fontenc}
  \usepackage[utf8]{inputenc}
\else
  \usepackage{fontspec}
\fi

\usepackage[margin=1in]{geometry}
\usepackage{booktabs}
\usepackage{longtable}
\usepackage{array}
\usepackage{amsmath}
\usepackage{microtype}
\usepackage{graphicx}
\usepackage{xcolor}
\usepackage{caption}
\usepackage[hidelinks]{hyperref}
\usepackage{url}

\captionsetup{font=small,labelfont=bf}
\setlength{\tabcolsep}{5pt}
\renewcommand{\arraystretch}{1.12}

% Sayı sütunları için sağa dayalı, sabit genişlikli hücre.
\newcolumntype{R}[1]{>{\raggedleft\arraybackslash}p{#1}}
\newcolumntype{L}[1]{>{\raggedright\arraybackslash}p{#1}}

% Kod ve sütun adları.
\newcommand{\col}[1]{\texttt{\small #1}}

\hypersetup{
  pdfauthor={Serdar I. Caglar},
  pdfkeywords={speech corpus, text-to-speech, Turkish, forced alignment, data curation}
}


\title{\textsc{Kiraat}: A 3,105-Hour Turkish Read-Speech Corpus\\
Segmented at Sentence Boundaries and Released Without Quality Gates}

\author{Serdar I. Ça\u{g}lar\\
\small Independent Researcher\\
\small \href{https://orcid.org/0000-0002-5776-2431}{ORCID: 0000-0002-5776-2431}}

\date{}

\begin{document}
\maketitle

\begin{abstract}
\noindent
Turkish has abundant text and speech-recognition data but very little read
speech suitable for training text-to-speech (TTS): the openly available
read-speech corpora for the language are measured in tens or low hundreds of
hours. We release \textsc{Kiraat}, a corpus of 1,840,404 clips totalling
3,105.7 hours of Turkish read speech, segmented from 2,680 audiobook and
spoken-article recordings across 27 channels, together with the pipeline that
produced it.

Two design rules distinguish the pipeline from the established recipe. First,
clips are cut at sentence boundaries recovered from word-level forced
alignment, not at silences found by voice activity detection. In a controlled
ablation over 206 recordings (240.5 hours) in which only the cut point varies,
14.31\% of the clips produced by silence-based segmentation begin mid-sentence
and 12.81\% do so with no signal that would let a consumer detect it; the
sentence-aligned arm leaves no unflagged broken clip at all. Over the full
corpus, 23 of the 1,547,494 clips in the recommended subset begin
mid-sentence, a rate of 15 per million.

Second, no threshold removes a clip. Every measurement stage writes columns
only, and the accept/exclude decision is made by a separate, versioned policy
read from configuration, so a consumer can reproduce a stricter or looser
subset without re-running a model. This separation is what let us test
thresholds instead of assuming them, and two standard ones did not survive:
in blind listening, a clipping-ratio gate rejected clips that listeners judged
usable 25 times out of 25, and an ASR-confidence gate rejected 15 of 30 clips
whose transcripts were verbatim correct. Both rules were removed from the
policy and both columns are still published.

We report the corpus, the measured limits of its source material, a
90-speaker clustering derived without fixing the number of speakers, an
evaluation split with zero recording and text leakage, and the pipeline cost
on a single GPU.
\end{abstract}

\section{Introduction}

Text-to-speech systems are trained on read speech: one speaker, a prepared
text, a quiet room or a good microphone, and sentences delivered whole. The
supply of such material is uneven across languages in a way that does not
track the number of speakers. Turkish is a case in point. It has large text
corpora, it is covered by the multilingual speech recognition and
representation models, and transcribed speech for recognition exists in
quantity; what it does not have is read speech in the form TTS training
needs. The openly available Turkish read-speech corpora are measured in tens
or low hundreds of hours, which is one to two orders of magnitude below what
comparable English resources offer.

The gap is therefore not in the language but in the task. This matters for how
one should try to close it. A method that only works where a language is
already well served is of little use; a method that turns publicly available
audiobook channels into a TTS-grade corpus, and that states its own defects in
measurable terms, transfers to any language that has such channels. That is
what we set out to build, and what we release here alongside the corpus
itself.

\paragraph{Contributions.}
\begin{enumerate}
\item \textbf{The corpus.} 1,840,404 clips, 3,105.7 hours of Turkish read
speech at 24\,kHz from 2,680 recordings on 27 channels, with 90 clustered
speaker identities, published with every per-clip measurement the pipeline
computed and with an evaluation split whose recording and text leakage are
both zero (Section~\ref{sec:corpus}).

\item \textbf{Segmentation at sentence boundaries, measured against the
alternative.} We reverse the usual order --- long-form ASR with word
timestamps first, sentence segmentation on the transcript second, cutting
third --- and quantify what that buys in a controlled ablation where the cut
point is the only variable (Section~\ref{sec:ablation}).

\item \textbf{A gate-free, versioned release policy.} Measurement stages are
structurally forbidden to make decisions; the recommended subset is produced
by a policy that lives in configuration, carries a version number, and records
a reason for every exclusion (Section~\ref{sec:policy}).

\item \textbf{Validated negative results.} Two thresholds in common use failed
blind listening on this corpus and were removed. We report the protocol, the
listener verdicts, and the columns that remain published so that others can
re-impose the thresholds if their own listening supports them
(Section~\ref{sec:negative}).

\item \textbf{A method for multi-hour recordings.} 30.7\% of our source
material is in single recordings longer than four hours, and the longest is
14.9 hours. Processing these under a fixed memory budget shapes the whole
architecture (Sections~\ref{sec:pipeline} and~\ref{sec:cost}).
\end{enumerate}

\section{Related work}

\paragraph{The established pipeline.}
Contemporary pipelines that turn found audio into speech corpora share a
skeleton: standardisation, optional source separation, speaker diarisation,
segmentation by voice activity detection, ASR, and quality filtering. Emilia
\cite{emilia,emilialarge} is the clearest statement of it, and GigaSpeech~2
\cite{gigaspeech2} instantiates the same order for languages with less data,
transcribing with Whisper \cite{whisper} and force-aligning with MMS
\cite{mms}. We share most of this skeleton and depart from it in two places.

\paragraph{Where to cut.}
Cutting at sentence boundaries rather than at silences is not new, and we do
not claim it as a contribution: LibriTTS \cite{libritts} separated itself from
LibriSpeech \cite{librispeech} on exactly this point, Libriheavy
\cite{libriheavy} places its cuts at sentence boundaries as well, and
WenetSpeech4TTS \cite{wenetspeech4tts} attacks the same defect from the other
side by merging adjacent segments and widening boundaries so that words are
not cut. What is missing from the literature, and what we supply, is a
controlled measurement of the cost of the alternative on the same recordings,
including the part that matters most to a consumer: whether the resulting
defect is detectable in the released metadata.

\paragraph{Filtering as a gate.}
Pipelines that build data to feed their own model filter inside the pipeline
with fixed thresholds. Emilia discards clips below a DNSMOS \cite{dnsmos} OVRL
of 3.0; LibriTTS requires an estimated SNR above 20\,dB for its ``clean''
subset, and its stricter path reduces LibriSpeech's 982 hours to 585. Projects
that release a shared corpus more often score and stratify instead:
WenetSpeech4TTS publishes Premium, Standard and Basic tiers cut at DNSMOS
P.808 thresholds and ships all three, and LibriTTS-R \cite{librittsr} takes
the same instinct further by restoring clips rather than discarding them. We
go one step beyond stratification: nothing is removed at any tier, the
thresholds live outside the data, and the policy that applies them is
versioned so that a published number can be reproduced exactly or
deliberately changed.

\paragraph{Turkish read speech.}
The Turkish speech resources in open circulation are dominated by
recognition-oriented material: read prompts collected by volunteers, broadcast
news, and multilingual evaluation sets. The read-speech portion suitable for
TTS training is small in every case, which is the gap this release addresses.
% TODO: her kaynağın saat sayısı birincil kaynaktan doğrulanıp buraya
% tablo olarak konacak (Common Voice tr, FLEURS tr, MediaSpeech tr,
% Turkish Speech Corpus). Doğrulanmadan sayı yazılmaz.

\section{Source material and its measured limits}
\label{sec:corpus}

The raw material is 2,700 files downloaded from 27 public Turkish audiobook
and spoken-article channels. Every file was probed before any processing
began: 2,698 are readable audio (one m4a is corrupt, one file is a README),
totalling \textbf{3,440.2 hours} and 200.1\,GB.

\paragraph{Recordings are long.}
The median recording is 41.4 minutes and the longest is 14.92 hours.
Table~\ref{tab:lengths} gives the distribution. The consequential figure is
that \textbf{30.7\% of the corpus hours sit in single recordings longer than
four hours}. Processing a multi-hour recording under a fixed memory budget is
therefore a requirement of this material, not an optimisation, and it shapes
the pipeline described in Section~\ref{sec:pipeline}. The established
pipelines do not help here: Emilia's inputs range from 20 seconds to just under
one hour, and Libriheavy resolves long-form audio by aligning it to the book
text, a reference we do not have.

\begin{table}[t]
\centering
\small
\begin{tabular}{lrrr}
\toprule
recording length & recordings & hours & share of corpus \\
\midrule
$<$ 10 min   & 168 & 20.4 & 0.6\% \\
10--30 min   & 911 & 274.0 & 8.0\% \\
30--60 min   & 635 & 462.8 & 13.5\% \\
1--2 h       & 467 & 654.7 & 19.0\% \\
2--4 h       & 351 & 973.4 & 28.3\% \\
4--8 h       & 131 & 700.5 & 20.4\% \\
8 h and over & 35  & 354.4 & 10.3\% \\
\bottomrule
\end{tabular}
\caption{Length distribution of the 2,698 readable source recordings.}
\label{tab:lengths}
\end{table}

\paragraph{The encoding profile is uniform, and that is a limitation.}
2,618 files are AAC (3,379.7 hours) and 80 are MP3 (60.5 hours); 2,693 are at
44.1\,kHz; the bit rate is tightly concentrated (5th percentile 128, median
129, 95th percentile 130\,kb/s). We measured the true bandwidth of the source
by taking a 60-second excerpt from the middle of one recording per channel and
computing a 32k-point FFT: the median cutoff frequency is \textbf{15.7\,kHz}
(minimum 13.0, maximum 15.8), and the median share of energy above 12\,kHz is
0.05\%. A 24\,kHz release rate (12\,kHz Nyquist) therefore sits below the
ceiling of the source everywhere in the corpus; upsampling further would buy
nothing and downsampling has no justification. The uniformity itself is a
limitation of the corpus and we state it in Section~\ref{sec:limits}: a model
trained on \textsc{Kiraat} sees one codec profile.

\paragraph{Channel shares are unequal, and no cap was applied.}
The largest channel contributes 19.5\% of the raw hours and the second 15.5\%;
the top two together contribute 35.2\%. We considered a per-channel hour cap
and decided against it: the cap that would have equalised the channels would
have reduced the corpus from 3,440 to 1,849 hours. Pruning channel diversity
with a threshold is precisely the move this pipeline exists to avoid, and the
imbalance is published rather than hidden --- every clip carries its
\col{channel} and \col{duration}, so a consumer can impose any cap they like.

\section{Pipeline}
\label{sec:pipeline}

The pipeline runs as a sequence of stages over a SQLite state store, with CPU
and GPU work overlapped by a scheduler that keeps one writer. Table~\ref{tab:stages}
lists the stages and the models, with the weight revisions pinned in the run
record.

\begin{table}[t]
\centering
\small
\begin{tabular}{L{2.1cm}L{5.1cm}L{7.4cm}}
\toprule
stage & what it does & model / setting \\
\midrule
\col{prepare} & decode, 24\,kHz mono, 40\,Hz high-pass, $-1$\,dB peak & ffmpeg \\
\col{asr} & long-form transcription with word timestamps & faster-whisper \texttt{large-v3}, rev.\ \texttt{edaa852e}, beam 5, temperature 0, \col{condition\_on\_previous\_text: false} \\
\col{align} & forced alignment, per-word score & torchaudio \texttt{MMS\_FA} (with \col{<star>}) \\
\col{boilerplate} & channel-level mining of station idents and announcements & word $n$-grams, threshold 40\% of a channel's recordings (min.\ 2), \col{max\_words: 24} \\
\col{segment} & cut at sentence boundaries, audio-based boundary refinement & \col{kiraat/segment.py}, \col{kiraat/boundaries.py} \\
\col{clip\_qc} & speech ratio, internal silence, edge silences, level, LUFS & silero VAD; BS.1770 \\
\col{music} & music-to-speech level and AudioSet music score & torchaudio \texttt{HDEMUCS\_HIGH\_MUSDB\_PLUS}; \texttt{MIT/ast-finetuned-audioset-10-10-0.4593} \\
\col{dnsmos} & P.835 SIG / BAK / OVRL & DNS-Challenge \texttt{sig\_bak\_ovr.onnx}, sha256 pinned in config \\
\col{speaker} & per-recording voice identity and within-recording consistency & ECAPA-TDNN \texttt{speechbrain/spkrec-ecapa-voxceleb}, rev.\ \texttt{5c0be387} \\
\col{export} & manifest and run record & \col{manifests/clips.jsonl}, \col{run.json} \\
\bottomrule
\end{tabular}
\caption{Pipeline stages. Every model revision is pinned and recorded in
\col{run.json} together with the git commit, the stage versions and the
configuration hash.}
\label{tab:stages}
\end{table}

\subsection{Segmentation at sentence boundaries}

The order is the point. A pipeline that segments first and transcribes second
can only cut where the audio is quiet, and quietness does not mark sentence
ends: a narrator breathes at a comma and may not breathe between two
sentences. We transcribe the whole recording first, obtain word-level
timestamps, segment the \emph{transcript} into sentences, and only then cut
the audio at the timestamps of the sentence boundaries. Voice activity
detection is retained, but only to find the speech regions handed to ASR.

Two situations force a cut that is not a sentence boundary. A sentence longer
than the duration ceiling is split at its internal punctuation, and the
resulting clips carry the \col{forced\_split} flag; a gap inside a sentence
that exceeds the tolerated silence forces a split flagged \col{gap\_split}. The
flags are the mechanism by which the defect stays visible:
Section~\ref{sec:integrity} shows that 97.6\% of all mid-sentence starts in
the corpus carry one of these two flags, and that the policy therefore keeps
them out of the recommended subset.

\subsection{Boundary refinement}

Cutting at a sentence boundary fixes \emph{which} point to cut at, not
\emph{exactly where} that point is: a forced aligner reports the end of a word
slightly early, so a cut placed on the timestamp clips the final phoneme. The
refinement stage looks at signal energy in a small window around the timestamp
between adjacent clips and pulls the cut to the actual silence, with the
window and thresholds in configuration.

This was found by ear, not by a metric. A first listening round found truncated
syllables; the refinement was written in response; a third round on 28 clips
found 28 clean --- zero broken starts, zero broken ends, zero truncated
syllables --- and a fourth round on 30 clips (20 suspected, 10 controls,
shuffled) found zero truncated words in all of them. The audit criterion moved
with the finding: it is no longer ``does the cut match the timestamp'' but
``does any word inside the clip's own word span fall outside the clip'', that
is, does the text match the audio.

\subsection{Duplicates}

A duplicate is not the same text read twice. Two narrators reading the same
sentence is prosodic diversity and is kept; a duplicate is the same text
\emph{and} the same audio. We define ``same audio'' by measurement --- duration,
LUFS, RMS and peak matching exactly --- and deliberately do not put the channel
in the key, since matching measurements mean the same recording regardless of
where it was published. Table~\ref{tab:dupe} shows why the definition had to be
measured rather than assumed: the obvious key over-counts by two orders of
magnitude, and a tolerant key collapses genuinely different readings.

\begin{table}[t]
\centering
\small
\begin{tabular}{lr}
\toprule
definition & clips counted as duplicates \\
\midrule
text + channel (the obvious key) & 72,839 \\
every measured value identical & 0 \\
audio measurements identical, within channel & 592 \\
audio measurements identical, channel ignored & 598 \\
duration rounded to 0.1\,s, LUFS to 0.1\,dB & 15,752 \\
\bottomrule
\end{tabular}
\caption{Duplicate counts over all 1,840,404 clips under five definitions. The
second row shows that a literal reading of the rule is unusable: a channel's
per-episode ident is byte-identical audio, but \col{word\_confidence} and
\col{leading\_silence\_sec} differ because they measure the clip's
surroundings, not its audio.}
\label{tab:dupe}
\end{table}

The released manifest flags \textbf{576} clips as duplicates, 453 of them in a
single channel's episode ident. 150 clips have no audio identity because their
audio could not be read, and are passed through without a duplicate check.

\section{Measurements, not decisions}
\label{sec:measurements}

Every clip stage writes columns and nothing else. This is enforced rather than
documented: the stage base class rejects any stage output containing a
\col{recommended} or \col{decision} field. Thresholds do not appear in code at
all; they live in a single configuration file whose loader refuses to silently
accept an unknown section or key. The column schema has one definition in the
source, and the dataset card is generated from it, so the published
documentation cannot drift from the published columns.

The consequence for a consumer is concrete. Every clip in the release carries
its speech ratio, internal and edge silences, peak, RMS and BS.1770 loudness,
clipping ratio, AudioSet music score and separated music-to-speech level,
DNSMOS SIG/BAK/OVRL, mean and minimum forced-alignment score, ASR word
confidence, word count, speaker identity, within-recording speaker
consistency, and its flags --- whether or not our policy recommends it.
Re-deriving a different subset requires re-reading a manifest, not re-running a
GPU stage.

\paragraph{Audio is measured, not modified.}
Source separation is run to \emph{measure} background music, and its output is
discarded; the released waveform is the decoded source at 24\,kHz, and loudness
is a published column rather than an applied normalisation. Pipelines that
normalise level and publish the separated signal make a choice on the
consumer's behalf that cannot be undone downstream; we make it undoable by not
making it.

\section{The recommended subset policy}
\label{sec:policy}

The policy is a versioned object in configuration. Version~6, the one used for
this release, has three rules: \col{speech\_ratio} $\ge$ 0.60,
\col{internal\_silence\_sec} $\le$ 1.0, and the absence of any of the flags
\col{oversize}, \col{forced\_split}, \col{gap\_split}, \col{short},
\col{duplicate}, \col{boilerplate}, \col{background\_music}. Applying it to the
full corpus yields the recommended subset: \textbf{1,547,494 clips, 2,575.2
hours}, which is 84.1\% of the clips and 82.9\% of the hours. Every excluded
clip carries the reason it was excluded.

Because the policy is separate from the data, its history is legible, and its
history is mostly a history of rules that were \emph{removed}. That is the
subject of the next section.

\section{Results}

\subsection{Segmentation: what the cut point costs}
\label{sec:ablation}

We ran both segmentation strategies over the same recordings, from the same
word timestamps, with the same duration settings; the only variable is where
the cut is placed. The sentence arm is the output of the full pipeline; the
baseline arm cuts the speech regions found by silero VAD at the silences
between them, splitting hard at the ceiling when a single region exceeds it.
The sample is channel-balanced: 206 recordings from all 27 channels, 240.5
hours. Table~\ref{tab:ablation} gives the result.

\begin{table}[t]
\centering
\small
\begin{tabular}{lrr}
\toprule
measure & sentence-aligned & silence-aligned \\
\midrule
clips & 141,379 & 76,466 \\
hours & 240.54 & 254.18 \\
begins mid-sentence & 1,967 (1.39\%) & 10,943 (\textbf{14.31\%}) \\
ends mid-sentence & 2,405 (1.70\%) & 12,117 (15.85\%) \\
\textbf{unflagged} broken start & \textbf{0 (0.00\%)} & 9,798 (\textbf{12.81\%}) \\
unflagged broken end & 0.01\% & 11.59\% \\
broken at both ends & 418 & 2,673 \\
median duration & 5.94\,s & 12.90\,s \\
duration, 5th--95th pct & 2.69--10.96\,s & 5.10--15.40\,s \\
at the duration ceiling & 0.64\% & 14.70\% \\
mean \col{align\_score} & 0.9442 & 0.9459 \\
mean of per-clip \emph{minimum} \col{align\_score} & \textbf{0.7319} & 0.6469 \\
\bottomrule
\end{tabular}
\caption{Segmentation ablation over 206 recordings (240.5 hours) from all 27
channels. Only the cut point differs between the two arms.}
\label{tab:ablation}
\end{table}

The row that matters is the unflagged one. In the sentence arm there is
\emph{no} unflagged clip with a broken start: all of the 1.39\% carry
\col{forced\_split} or a boilerplate flag, so the defect is visible in the
release and excluded from the recommended subset. In the baseline arm almost
all of the breakage is unflagged (12.81\%), because a pipeline that cuts at
silence has no signal with which to mark it. The defect varies by channel but
never disappears: the baseline arm's mid-sentence rate ranges from 7.0\% to
35.0\% across channels, the sentence arm's from 0.00\% to 5.02\%.

Two side findings. Silence-aligned segmentation yields \emph{more} hours
(254.2 against 240.5) from roughly half as many clips: its clips are twice as
long and 14.7\% of them end at the duration ceiling, meaning the cut point is
set by the ceiling rather than by the content. And while mean alignment
confidence is identical across the two arms, the mean of the per-clip
\emph{minimum} word confidence is clearly higher in the sentence arm (0.732
against 0.647): clips cut at silences carry more badly aligned words at their
edges.

\subsection{Sentence integrity over the full corpus}
\label{sec:integrity}

The ablation is a controlled comparison on a balanced sample. We also applied
the same two definitions --- a clip whose text begins with a lowercase letter,
and a clip whose text does not end in sentence-final punctuation, both computed
with the same Turkish-aware routines --- to all 1,840,404 clips.

\begin{table}[t]
\centering
\small
\begin{tabular}{lrrrr}
\toprule
subset & clips & hours & broken start & broken end \\
\midrule
whole corpus & 1,840,404 & 3,105.7 & 27,483 (1.49\%) & 33,346 (1.81\%) \\
recommended subset & 1,547,494 & 2,575.2 & \textbf{23 (0.0015\%)} & 430 (0.03\%) \\
outside the policy & 292,910 & 530.5 & 27,460 (9.37\%) & 32,916 (11.24\%) \\
\bottomrule
\end{tabular}
\caption{Sentence integrity over the full corpus. Rates are per clip.}
\label{tab:integrity}
\end{table}

Table~\ref{tab:integrity} should be read together with the flags.
Of the 27,483 clips that begin mid-sentence, 97.6\% carry \col{forced\_split}
or \col{gap\_split} (92.3\% carry \col{forced\_split} alone); among the
1,558,409 clips that carry no flag at all, 23 begin mid-sentence, a rate of
15 per million. The corpus does contain
mid-sentence cuts --- a sentence that runs past the duration ceiling has to be
split somewhere --- but they are cuts the pipeline declares, and declaring them
is what allows the policy to exclude them and a consumer to find them.

\subsection{Blind listening: two standard thresholds that did not hold}
\label{sec:negative}

Because the thresholds are separate from the data, they can be tested. Our
protocol hides the channel and the measurement, balances the clips across
bands of the measurement under test, and records the listener's verdict to a
file before the scores are revealed. Six rounds were run: two on background
music, two on boundary cuts, one on clipping ratio and one on ASR word
confidence. Two of them ended with a rule being deleted.

\paragraph{Clipping ratio.}
A rule excluding clips with a clipping ratio above 0.002 was removing 536
clips --- and 532 of them were in a single channel, 93.2\% of that channel's
output. That is the signature of a channel property, not a clip defect. A blind
round of 25 clips balanced across six bands of the measure (0 to 0.0227) found
24 ``no audible distortion'' and one ``slight'', and that one was in the lowest
band; all 25 were judged fit for training. The rule was removed in policy v5.

\paragraph{ASR word confidence.}
A rule requiring word confidence $\ge$ 0.60 was tested the same way: 30 clips
across six bands from 0.18 to 0.99, channel and score hidden. In 29 the
transcript was verbatim correct; the single error was in the lowest-confidence
clip. The threshold would have removed 15 of these 30, of which 14 were found
flawless. An independent check without ears agrees that the signal is weak but
real: rank correlation with the aligner score is $+0.262$, and even in the
worst band the median \col{align\_score\_mean} is 0.938. The rule was removed in
policy v6, and the column is still published as a ranking signal.

\paragraph{What survived.}
The background-music rule survived the same treatment and was corrected by it.
Our first guess of $-30$\,dB was refuted: in the $-40$ to $-33$\,dB band, 6 of
6 clips had clearly audible music. A second round of 36 clips balanced across
four level bands found audible music in 86\% of the flagged clips and clean
audio in 11 of 14 unflagged ones. The threshold stayed at $-40$\,dB, and the
rule is the conjunction of two columns --- separated music-to-speech level
above $-40$\,dB \emph{and} AudioSet music score $\ge$ 0.3 --- because the two
columns measure different things: their Spearman correlation on 3,097
separated clips is only $+0.479$.

\paragraph{A rule we did not re-impose.}
The DNSMOS OVRL $\ge$ 3.0 gate used by other pipelines is not in our policy.
The column exists for every clip and 83.57\% of the corpus would pass it; we
report that number precisely so that the gate's effect is visible, and we will
not adopt it until it passes the same blind listening the other rules were
subjected to. The burden of proof is on the gate, not on the data.

\subsection{The released columns}

Table~\ref{tab:cols} gives the distribution of the main columns over the
recommended subset. Two of them carry design decisions. Loudness is published
and \emph{not} applied: the median clip is at $-20.92$\,LUFS with a 5th--95th
percentile range of $-27.99$ to $-12.54$, and normalising is the consumer's
choice. DNSMOS is published for every clip: over the whole corpus the mean OVRL
is 3.208 (5th percentile 2.69, median 3.27, 95th 3.49).

\begin{table}[t]
\centering
\small
\begin{tabular}{lrrrr}
\toprule
column & p05 & p50 & p95 & mean \\
\midrule
duration (s) & 2.76 & 5.82 & 10.56 & 5.99 \\
\col{loudness\_lufs} & $-27.99$ & $-20.92$ & $-12.54$ & $-20.61$ \\
\col{dnsmos\_ovrl} & 2.862 & 3.296 & 3.491 & 3.253 \\
\col{dnsmos\_sig} & 3.271 & 3.563 & 3.710 & 3.536 \\
\col{dnsmos\_bak} & 3.631 & 4.109 & 4.210 & 4.041 \\
\col{speech\_ratio} & 0.851 & 0.977 & 1.000 & 0.955 \\
\col{music\_score\_audioset} & 0.001 & 0.006 & 0.228 & 0.034 \\
\col{music\_to\_speech\_db} & $-80.0$ & $-80.0$ & $-40.94$ & $-75.42$ \\
\col{align\_score\_mean} & 0.858 & 0.961 & 0.994 & 0.946 \\
\col{align\_score\_min} & 0.278 & 0.800 & 0.976 & 0.740 \\
\col{word\_confidence} & 0.550 & 0.917 & 0.999 & 0.866 \\
\col{n\_words} & 5 & 11 & 21 & 11.7 \\
\col{internal\_silence\_sec} & 0.000 & 0.000 & 0.664 & 0.166 \\
\bottomrule
\end{tabular}
\caption{Column distributions over the recommended subset (1,547,494 clips).}
\label{tab:cols}
\end{table}

The flags over the whole corpus are: \col{background\_music} 207,196,
\col{forced\_split} 52,501, \col{short} 20,904, \col{oversize} 6,731,
\col{gap\_split} 5,141, \col{boilerplate} 1,053, \col{duplicate} 576. The
music flag alone accounts for most of the difference between the corpus and
the recommended subset.

\subsection{Speakers}

Channels are not speakers, and we did not assume they were. The \col{speaker}
stage computes an ECAPA-TDNN centroid per recording and a within-recording
consistency; the identity decision is made afterwards, at corpus level, by
average-linkage clustering. The merge threshold is not chosen by hand and the
number of clusters is never supplied: we measure the similarity distribution
within recordings and across recordings, and take the threshold at the equal
error point of the two.

On the full corpus (2,677 recordings; 21 have no clip long enough to establish
an identity) the equal error point is at \textbf{0.532} with an error of
0.52\%, and clustering at that threshold yields \textbf{90 speakers}. The
median margin is 0.445 and no recording has a negative margin.

Two findings bear on the design. \textbf{13 of the 27 channels contain more
than one speaker}, and \textbf{9 clusters span several channels} --- the same
narrator publishes on two or three of them. Channel balance is therefore not
speaker balance, and neither is a substitute for the other. The distribution is
long-tailed: the largest speaker holds 18.2\% of the hours, the top five 54.4\%
and the top ten 75.0\%, while 17 speakers have less than an hour.

\subsection{Evaluation split}

The split is at the recording level: all clips of a recording go to one split,
so evaluation is on a recording the model has never heard. The \col{dev} and
\col{test} splits are channel-balanced, because the corpus is not and an
evaluation set that inherited the imbalance would report the numbers of two
channels. Recordings run for hours, so dev and test are reduced to a per-channel
hour target and the clips beyond that target are left out of training rather
than moved into it; they are published in the \col{rest} split described
below.

\begin{table}[t]
\centering
\small
\begin{tabular}{lrrrr}
\toprule
split & clips & hours & recordings & channels \\
\midrule
train & 1,512,448 & 2,517.08 & 2,579 & 27 \\
dev & 5,750 & 9.54 & 40 & 27 \\
test & 5,473 & 9.30 & 45 & 27 \\
\midrule
rest & 316,733 & 569.78 & --- & 27 \\
\textbf{total} & \textbf{1,840,404} & \textbf{3,105.70} & \textbf{2,680} & \textbf{27} \\
\bottomrule
\end{tabular}
\caption{The four published splits. \col{train}, \col{dev} and \col{test} are
drawn from the recommended subset; \col{rest} carries every remaining clip, so
the four splits together are the whole corpus. The recordings column counts
recordings whose clips form a split; \col{rest} draws clips from recordings
across all three, and 16 of the 2,680 recordings have all of their clips in it.}
\label{tab:splits}
\end{table}

Nothing is left out of the release. The clips that the policy does not
recommend, and the recommended clips that dev and test did not use, are
published in a fourth split, \col{rest}: 292,910 clips (530.5 hours) with
\col{recommended=false}, and 23,823 clips (39.3 hours) that are recommended but
fell beyond a channel's dev/test duration target or were dropped by the
text-leakage cleanup below. A clip in \col{rest} carries exactly the same
measurement columns as one in \col{train}, so a consumer who sets a different
policy can recover it from the release rather than re-running the pipeline ---
which is the whole point of not gating.

Recording leakage is zero by construction and is nevertheless tested. Text
leakage was not zero by construction: 2,263 test clips and 2,723 dev clips
carried a sentence that appears verbatim in training, because the same book is
read by different narrators on different channels. They were dropped from the
evaluation sets --- training was left untouched --- and the remaining leakage is
zero.

Speaker overlap is a different matter and we report it plainly: train has 89
speakers, dev 26 and test 26, and \textbf{25 of the 26 test speakers also occur
in training}. This is the usual setup for TTS evaluation, but it is a
\emph{seen-speaker} evaluation and must not be described otherwise. An
unseen-speaker evaluation would require a separately constructed subset.

\section{Cost and reproducibility}
\label{sec:cost}

The full corpus was processed on one machine: an RTX 3090 (24\,GB) with 12 CPU
workers, under a 280\,W power cap. Forced alignment is the bottleneck at 94
audio-hours per hour with a single worker. Profiling it showed 56\% of the time
in the model forward pass, 28\% in token merging and 10\% in the alignment
itself; the token-merging cost was a GPU synchronisation per token span, and
moving the score tensor to CPU raised the same recording from 109$\times$ to
188$\times$ real time with byte-identical output. The other stages measured
46 clips/s for clip QC across 12 CPU workers, 28 clips/s for music, 81.1
clips/s for DNSMOS and about one second per recording for the speaker stage.

The working directory occupies 567\,GB: 284\,GB of intermediate 24\,kHz mono
FLAC, 246\,GB of clips (about 159\,KB per clip), 2.7\,GB of alignment, 1.5\,GB
of ASR output and a 1.4\,GB state database.

Every run writes a record containing the git commit, the version of each stage,
the configuration hash and the pinned revision of every model weight, so that a
published number can be traced to the code and the weights that produced it.

\section{Limitations}
\label{sec:limits}

\paragraph{No downstream TTS experiment yet.}
This paper describes and measures a corpus and the pipeline that built it. It
does not yet report a TTS system trained on that corpus. The comparison that
would close the argument --- equal-budget models trained on the recommended
subset and on the whole corpus, evaluated with MOS/CMOS, intelligibility and
speaker similarity --- is designed but not run, and until it is, the claims here
are claims about data, not about synthesis quality.

\paragraph{Numerals are not force-aligned.}
The aligner's lexicon has no digits, so words containing numerals fall back to
the ASR timestamp. Rewriting them into their spoken form before alignment was
rejected because it would alter the published text.

\paragraph{Alignment noise at sentence starts.}
245 words in 289 clips (2.12\% of a probed sample) receive implausible
alignments on short sentence-initial words. The cause is not established; a
chunk-boundary hypothesis was tested and rejected.

\paragraph{Channel and speaker concentration.}
The top two channels are 38.2\% of the recommended subset and the top five
speakers are 54.4\% of the hours. Both are published per clip so that a
consumer can rebalance, but the corpus as released is concentrated.

\paragraph{One encoding profile.}
The entire corpus is roughly 129\,kb/s AAC with a median source bandwidth of
15.7\,kHz, and two channels are band-limited near 13\,kHz.

\paragraph{Seen-speaker evaluation.}
25 of the 26 test speakers appear in training, as reported above.

\paragraph{Boilerplate detection is n-gram based.}
Verbatim $n$-gram mining does not catch templated announcements with variable
slots (``from the book \textit{X} by \textit{Y}''). 23 clips in the recommended
subset still begin mid-sentence, and they come from this path.

\section{Conclusion}

\textsc{Kiraat} is 3,105.7 hours of Turkish read speech, of which 2,575.2 hours
are recommended by a policy that anyone can re-run, tighten or discard. The two
rules that shaped it are simple to state and were expensive to establish: cut
where sentences end, and let no threshold delete a clip. The first is
measurable --- on the same recordings, silence-based segmentation leaves 12.81\%
of clips broken and undetectable, against none in the sentence-aligned arm. The
second is what allowed the thresholds themselves to be tested, and two standard
ones did not survive contact with a listener.

\section*{Availability}

% TODO: HF veri kümesi bağlantısı, kod deposu bağlantısı, lisans satırı ve
% kanal künyesi şartı yayın anında buraya yazılacak.

\bibliographystyle{plain}
\bibliography{refs}

\end{document}
